\scanpage{107}

\mspage{113}

Je rappelle aussi la

\uline{Définition} 1 : Soit $f:\deleted{Z}X\to Y$ de $\mathrm{Diag}$. On dit que $f$ est une $\mathbb D$-\uline{équiv. cohomologique} (sous-entendu : ``\uline{complète}'') si pour tout $\mathbb D$-coeff. $F$ \uline{loc constant} sur $Y$,

\[
\deleted{\unsure{F|K}}\quad H^{\unsure{\cdot|i}}_{\mathbb D}(Y,F)\xrightarrow{\sim}H^{\cdot}_{\mathbb D}(X,f^{\unsure{\cdot|{\prime}|*}}F)\ ;
\]

on dit que $f$ est une $\mathbb D$-\uline{équivalence homologique} (``\uline{complète}'') si pour tout $F$ comme dessus,

\[
H_{\cdot}^{\mathbb D}(X,f^{\cdot}F)\xrightarrow{\sim}H_{\cdot}^{\mathbb D}(Y,F).
\]

On dit que $f$ est une $\mathbb D$-\uline{équivalence $1$-complète} (resp. $2$-\uline{complète}) si \deleted{pour} le foncteur

\[
f^*:\mathbb D(Y)_{\mathrm{loc.const.}}\longrightarrow\mathbb D(X)_{\mathrm{loc.const}}
\]

est pleinement fidèle (resp. une équivalence).

\diagram{\begin{tikzcd}[ampersand replacement=\&,row sep=small,column sep=small]
 \& \mathbb D\text{-équiv. }\unsure{2|3}\text{-complète}\arrow[d,Rightarrow] \& {}\\
 \& \mathbb D\text{-équiv. }1\text{-complète}\arrow[dl,Rightarrow]\arrow[dr,Rightarrow] \& {}\\
 \mathbb D\text{-équiv. coh.}\arrow[dr,Rightarrow] \& \& \mathbb D\text{-équiv. hom.}\arrow[dl,Rightarrow]\\
 \& \mathbb D\text{-équiv.} \&
\end{tikzcd}}

Si dans les \inserted{sous-catégories} $\mathbb D(X)_{\mathrm{loc.const}}$ il y a une loi $*$ (associative, unitaire, commutative) et des $\mathcal H\mathrm{om}$ internes associés, préservés par les images inverses, alors les $\mathbb D$-équiv. \deleted{$1$-complètes} cohomologiques sont déjà $1$-complètes, et sont \inserted{\unsure{bien|donc|même|aussi}} des $\mathbb D$-équiv. homologiques, donc dans ce cas particulier

% archive sheet number: 107
